% WHAT LIST A DISPLAY STANDS BESIDE.
%
% A display ends the paragraph it interrupted and stands BESIDE its lines
% in the vertical list that holds them -- so the list a display is in is
% the one the paragraph was in, INTERNAL where the paragraph was inside a
% \vbox. \tracingcommands says so: the `}' that ends
%
%   \vbox{\noindent A$$B$$\ifvmode\fi C\par}
%
% draws `{internal vertical mode: end-group character }}', and the blank
% space after it draws `{vertical mode: blank space  }' because the mode
% has changed.
%
% This engine let the paragraph's own mode stand while the display took
% over, so the level holding the display was left saying `vertical mode'
% and everything after it in that \vbox was named wrongly. trip line 397
% writes a display inside a \vbox and asks \ifvmode right after it.
\catcode`\{=1 \catcode`\}=2 \catcode`\$=3
\tracingonline=1
\font\r=cmr10 \font\s=cmsy10 \font\x=cmex10 \r
\textfont0=\r \scriptfont0=\r \scriptscriptfont0=\r
\textfont1=\r \scriptfont1=\r \scriptscriptfont1=\r
\textfont2=\s \scriptfont2=\s \scriptscriptfont2=\s
\textfont3=\x \scriptfont3=\x \scriptscriptfont3=\x
\hsize=100pt \parindent=0pt \hbadness=10000 \vbadness=10000
\message{[A a display inside a vbox]}
\tracingcommands=2
\setbox0=\vbox{\noindent A$$B$$\ifvmode\fi C\par}
\tracingcommands=0
\message{[B one at the outer level]}
\vsize=20000pt \output={\shipout\box255}
\tracingcommands=2
\noindent A$$B$$\ifvmode\fi C\par
\tracingcommands=0
\message{[done]}
\end
